% GENERATED FILE. Edit canonical lessons, then run npm run generate:books.
% canonical-source-hash: fnv1a64:a3691f43ce95c7ba
% canonical-lessons: PT-C168-polvo, PT-C168-alface, PT-C168-espinafres, PT-C168-abobora, PT-C168-ervilha

\chapter{Octopus, Lettuce, Spinach, Pumpkin, Pea}
\label{ch:pt-octopus-lettuce-spinach-pumpkin-pea}
\hlchaptermodality{168}

\begin{chapteropening}
\textbf{By the end of this chapter:} \emph{I can say an octopus, lettuce, spinach, a pumpkin and a pea.}

Complete the last lesson of chapter 168: I can say an octopus, lettuce, spinach, a pumpkin and a pea.
\end{chapteropening}

\section[o polvo]{o polvo — an octopus}

\label{lesson:PT-C168-polvo}



\begin{quote}
Before the new one: say the Portuguese for a brand, then the Portuguese for quality.
\end{quote}

\subsection*{You'll want to know: o polvo}
\textbf{o polvo} — \textquotedblleft{}an octopus\textquotedblright{}.

\begin{cousinweb}
\textbf{polvo} $\leftarrow$ Latin \emph{polypus}, \textquotedblleft{}octopus\textquotedblright{} (EtymDB 2.1) — inherited. Latin took it from Greek \emph{polypous}, \textquotedblleft{}many-footed\textquotedblright{}. Old Portuguese said \emph{polbo}.
\end{cousinweb}

\begin{grammarlens}[title={What you've built}]
One more word for this chapter.
\end{grammarlens}

\subsection*{Your turn}
\begin{itemize}
\raggedright
  \item \emph{Say it:} \emph{o polvo}
  \item \emph{Say it:} \emph{o polvo}, once more
  \item \emph{Say it:} say \emph{o polvo}
  \item \emph{Recall:} say the Portuguese for a brand, then the Portuguese for quality, then say \emph{o polvo} again
\end{itemize}

\begin{tcolorbox}[breakable,colback=teal!4,colframe=teal!35!black,title={Before you move on}]
What does o polvo mean? (\textquotedblleft{}An octopus\textquotedblright{}.) Say it once more.
\end{tcolorbox}
\section[a alface]{a alface — lettuce}

\label{lesson:PT-C168-alface}



\begin{quote}
Before the new one: say the Portuguese for quality, then the Portuguese for an octopus.
\end{quote}

\subsection*{You'll want to know: a alface}
\textbf{a alface} — \textquotedblleft{}lettuce\textquotedblright{}.

\begin{cousinweb}
\textbf{alface} $\leftarrow$ Arabic \emph{al-khass}, \textquotedblleft{}the lettuce\textquotedblright{} (EtymDB 2.1) — borrowed. The \emph{al-} is the Arabic word for \textquotedblleft{}the\textquotedblright{}.
\end{cousinweb}

\begin{grammarlens}[title={What you've built}]
One more word for this chapter.
\end{grammarlens}

\subsection*{Your turn}
\begin{itemize}
\raggedright
  \item \emph{Say it:} \emph{a alface}
  \item \emph{Say it:} \emph{a alface}, once more
  \item \emph{Say it:} say \emph{a alface}
  \item \emph{Recall:} say the Portuguese for quality, then the Portuguese for an octopus, then say \emph{a alface} again
\end{itemize}

\begin{tcolorbox}[breakable,colback=teal!4,colframe=teal!35!black,title={Before you move on}]
What does a alface mean? (\textquotedblleft{}Lettuce\textquotedblright{}.) Say it once more.
\end{tcolorbox}
\section[os espinafres]{os espinafres — spinach}

\label{lesson:PT-C168-espinafres}



\begin{quote}
Before the new one: say the Portuguese for an octopus, then the Portuguese for lettuce.
\end{quote}

\subsection*{You'll want to know: os espinafres}
\textbf{os espinafres} — \textquotedblleft{}spinach\textquotedblright{}.

\begin{cousinweb}
\textbf{espinafres} is the plural of the spinach word $\leftarrow$ Arabic \emph{isfānākh} — borrowed; Arabic took it from Persian.
\end{cousinweb}

\begin{grammarlens}[title={What you've built}]
One more word for this chapter.
\end{grammarlens}

\subsection*{Your turn}
\begin{itemize}
\raggedright
  \item \emph{Say it:} \emph{os espinafres}
  \item \emph{Say it:} \emph{os espinafres}, once more
  \item \emph{Say it:} say \emph{os espinafres}
  \item \emph{Recall:} say the Portuguese for an octopus, then the Portuguese for lettuce, then say \emph{os espinafres} again
\end{itemize}

\begin{tcolorbox}[breakable,colback=teal!4,colframe=teal!35!black,title={Before you move on}]
What does os espinafres mean? (\textquotedblleft{}Spinach\textquotedblright{}.) Say it once more.
\end{tcolorbox}
\section[a abóbora]{a abóbora — a pumpkin}

\label{lesson:PT-C168-abobora}



\begin{quote}
Before the new one: say the Portuguese for lettuce, then the Portuguese for spinach.
\end{quote}

\subsection*{You'll want to know: a abóbora}
\textbf{a abóbora} — \textquotedblleft{}a pumpkin\textquotedblright{}.

\begin{cousinweb}
\textbf{abóbora} — its origin is uncertain. EtymDB 2.1 has a row for it but no ancestor.
\end{cousinweb}

\begin{grammarlens}[title={What you've built}]
One more word for this chapter.
\end{grammarlens}

\subsection*{Your turn}
\begin{itemize}
\raggedright
  \item \emph{Say it:} \emph{a abóbora}
  \item \emph{Say it:} \emph{a abóbora}, once more
  \item \emph{Say it:} say \emph{a abóbora}
  \item \emph{Recall:} say the Portuguese for lettuce, then the Portuguese for spinach, then say \emph{a abóbora} again
\end{itemize}

\begin{tcolorbox}[breakable,colback=teal!4,colframe=teal!35!black,title={Before you move on}]
What does a abóbora mean? (\textquotedblleft{}A pumpkin\textquotedblright{}.) Say it once more.
\end{tcolorbox}
\section[a ervilha]{a ervilha — a pea}

\label{lesson:PT-C168-ervilha}



\begin{quote}
Before the new one: say the Portuguese for spinach, then the Portuguese for a pumpkin.
\end{quote}

\subsection*{You'll want to know: a ervilha}
\textbf{a ervilha} — \textquotedblleft{}a pea\textquotedblright{}.

\begin{cousinweb}
\textbf{ervilha} $\leftarrow$ Latin \emph{ervilia}, \textquotedblleft{}a kind of vetch\textquotedblright{} (EtymDB 2.1) — inherited. The \emph{-li-} became the \emph{lh} sound.
\end{cousinweb}

\begin{grammarlens}[title={What you've built}]
One more word for this chapter.
\end{grammarlens}

\subsection*{Your turn}
\begin{itemize}
\raggedright
  \item \emph{Say it:} \emph{a ervilha}
  \item \emph{Say it:} \emph{a ervilha}, once more
  \item \emph{Say it:} say \emph{a ervilha}
  \item \emph{Recall:} say the Portuguese for spinach, then the Portuguese for a pumpkin, then say \emph{a ervilha} again
\end{itemize}

\begin{tcolorbox}[breakable,colback=teal!4,colframe=teal!35!black,title={Before you move on}]
What does a ervilha mean? (\textquotedblleft{}A pea\textquotedblright{}.) Say it once more.
\end{tcolorbox}
